\subsection{连续映射的紧集}

定理 2（阿斯科利）. 设 X 是拓扑（相应地，一致）空间，S 是 X 的一个覆盖，Y 是一致空间，H 是 X 到 Y 中的映射所成之集，使得对每个 A ∈ S 及每个 u ∈ H，u 在 A 上的限制连续（相应地，一致连续）。则为使 H 关于 S-收敛的一致结构为预紧，下列条件在任何情形都是必要的，且当诸集 \(\mathbf{A} \in \mathfrak{S}\) 紧（相应地，预紧）时也是充分的：

\begin{enumerate}
\def\labelenumi{\alph{enumi})}
\item
  对每个 \(\mathbf{A} \in \mathfrak{S}\) ，集 \(\mathbf{H}|\mathbf{A} \subset \mathcal{F}(\mathbf{A};\mathbf{Y})\) （在 \(\mathbf{A}\) 上的、\(\mathbf{H}\) 中诸函数的限制）是等度连续（相应地，一致等度连续）的。
\item
  对每个 \(x \in \mathbf{X}\) ，集 \(\mathbf{H}(x) \subset \mathbf{Y}\) （由点 \(u(x)\) （ \(u \in \mathbf{H}\) ）组成）是预紧的。
\end{enumerate}

\begin{enumerate}
\def\labelenumi{\arabic{enumi})}
\tightlist
\item
  先证条件 a）与 b）是必要的。我们已知（§ 1，第 2 号，注 6）映射 \(u \to u(x)\) （由 \(\mathcal{F}_{\mathfrak{S}}(\mathbf{X};\mathbf{Y})\) 到 Y 中）是一致连续的；因此，若 H 预紧，则 \(\mathbf{H}(x)\) 也预紧（第 II 章，§ 4，第 2 号，命题 2），这就证明了 b）。为证 a），考虑一个集 \(\mathbf{A} \in \mathfrak{S}\) 、一点 \(x_0 \in \mathbf{A}\) 以及 Y 的一个围域 V；由于 H 预紧，它可以被有限个 \(W(\mathbf{A},\mathbf{V})\) -小的集所覆盖；换句话说，存在 H 中元素的一个有限序列 \((u_i)\) ，使得对每个 \(u \in \mathbf{H}\) ，我们有
\end{enumerate}

\[
(u (x), u _ {i} (x)) \in \mathbf {V} \quad \text { for   all } \quad x \in \mathbf {A}\tag{4}
\]

对至少一个指标 i 成立。

由于每个 \(u_{i}|A\) 都在 \(x_{0}\) 处连续（相应地，一致连续），存在邻域 \(U_{i}\) （\(x_{0}\) 在 A 中的邻域）（相应地，A 的一个围域 \(M_{i}\) ），使得

\[
x \in \mathrm{U} _ {i} \quad \text { implies } \quad (u _ {i} (x), u _ {i} (x _ {0})) \in \mathrm{V},\tag{5}
\]

（相应地，使得

\[
(x ^ {\prime}, x ^ {\prime \prime}) \in \mathbf {M} _ {i} \quad \text { implies } \quad (u _ {i} (x ^ {\prime}), u _ {i} (x ^ {\prime \prime})) \in \mathbf {V}.\tag{6}
\]

设 U（相应地，M）是诸 \(U_{i}\) （相应地，诸 \(M_{i}\) ）的交；它是 \(x_{0}\) 在 A 中的一个邻域（相应地，A 的一个围域）。对每个 \(u \in H\) 存在一个使（4）成立的指标 i；把条件（4）对 \(x_{0}\) 以及对 x（相应地，对 \(x'\) 与 \(x''\) ）写出，并注意到（5）{[}相应地（6）{]}，我们立即看出，关系 \(x \in U\) {[}相应地 \((x', x'') \in M\) {]}蕴涵 \((u(x), u(x_{0})) \in \overset{\mathbf{3}}{\mathrm{V}}\) {[}相应地 \((u(x'), u(x'')) \in \overset{\mathbf{3}}{\mathrm{V}}\) {]}对每个 \(u \in H\) 成立；这就证明了 a）。

\begin{enumerate}
\def\labelenumi{\arabic{enumi})}
\setcounter{enumi}{1}
\tightlist
\item
  现在证明：当诸集 \(A \in S\) 紧（相应地，预紧）时，条件 a）与 b）是充分的。条件 b）蕴涵 H 关于逐点收敛的一致结构是预紧的（第 II 章，§ 4，第 2 号，命题 3）。但由条件 a）与第 4 号的定理 1 可知，在 H\textbar A 上，A 中逐点收敛的一致结构与 A 中一致收敛的一致结构相同；因此 H\textbar A 在 \(\mathcal{F}_{u}(A; Y)\) 中是预紧的，这蕴涵 H 关于 S-收敛的一致结构是预紧的（§ 1，第 2 号）。
\end{enumerate}

注意，若 Y 是预紧空间，则定理 2 的条件 b）自动满足。

推论 1. 设 X 是拓扑（相应地，一致）空间，Y 是豪斯多夫一致空间，H 是 \(\mathcal{C}(\mathbf{X};\mathbf{Y})\) 的一个等度连续（相应地，一致等度连续）子集。假设 \(\mathrm{H}(x)\) 对每个 \(x\in \mathbf{X}\) 在 Y 中相对紧。则 H 在 \(\mathcal{C}(\mathbf{X};\mathbf{Y})\) 中关于紧（相应地，预紧）收敛拓扑是相对紧的。

设 \(\overline{\mathbf{H}}\) 是 \(\mathbf{H}\) 在 \(\mathcal{F}_s(\mathbf{X};\mathbf{Y})\) 中的闭包。\(\overline{\mathbf{H}}\) 是等度连续（相应地，一致等度连续）的（第 3 号，命题 6）。此外，我们有 \(\overline{\mathbf{H}}(x)\subset \overline{\mathbf{H}(x)}\) （ \(\S\) 1，第 2 号，注 6），因此 \(\overline{\mathbf{H}}(x)\) 也相对紧；于是定理 2 表明 \(\overline{\mathbf{H}}\) 关于 \(\mathfrak{S}\) -收敛是预紧的，这里 \(\mathfrak{S}\) 表示 \(\mathbf{X}\) 的所有紧（相应地，预紧）子集所成之集。此外，由于 \(\overline{\mathbf{H}(x)}\) 是紧的从而是完备的，\(\overline{\mathbf{H}}\) 关于逐点收敛的一致结构是完备的（第 II 章，\(\S\) 3，第 5 号，命题 10 以及第 4 号，命题 8），从而关于 \(\mathfrak{S}\) -收敛的一致结构也是完备的（ \(\S\) 1，第 5 号，命题 5 的推论 2）；于是 \(\overline{\mathbf{H}}\) 是紧的，因为它是预紧、完备且豪斯多夫的（ \(\S\) 1，第 2 号，命题 1）。

推论 2. 设 X 是拓扑（相应地，一致）空间，Y 是完备豪斯多夫一致空间，H 是 C(X; Y) 的一个等度连续（相应地，一致等度连续）子集。假设对一切 \(x \in D\) ，H(x) 在 Y 中相对紧，其中 D 是 X 的一个稠密子集。则 H 在 C(X; Y) 中关于紧（相应地，预紧）收敛拓扑是相对紧的。

只需证明 \(\mathrm{H}(x)\) 对一切 \(x \in X\) 相对紧，因为这样就可应用推论 1。由于 Y 完备，只需证明 \(\mathrm{H}(x)\) 对一切 \(x \in X\) 预紧。设 V 是 Y 的任一对称围域，则存在 x 的邻域 U，使得 \((u(x), u(x')) \in V\) 对一切 \(x' \in U\) 及一切 \(u \in H\) 成立。由假设存在 \(x' \in U \cap D\) ，且由于 \(\mathrm{H}(x')\) 在 Y 中相对紧，存在有限个点 \(y_k \in Y\) ，使得 \(\mathrm{H}(x')\) 含于诸集 \(\mathrm{V}(y_k)\) 的并；于是 \(\mathrm{H}(x)\) 含于诸集 \(\overset{2}{\mathrm{V}}(y_k)\) 的并，证明完毕。

推论 3. 设 X 是局部紧空间，Y 是豪斯多夫一致空间，H 是 C (X; Y) 的一个子集。则 H 在 \(\mathcal{C}_c(X; Y)\) 中相对紧当且仅当 H 等度连续且 H(x) 对一切 \(x \in \mathbf{X}\) 在 Y 中相对紧。

鉴于推论 1，只需证明：若 H 在 \(\mathcal{C}_{c}(\mathbf{X};\mathbf{Y})\) 中相对紧，则 H 等度连续。现在每一点 \(x\in X\) 都有一个紧邻域 A，且由定理 2 知 H/A 等度连续；这蕴涵 H 在 x 处等度连续，结论得证。

注。设 X 是拓扑空间，Y 是一致空间，S 是 X 的子集所成之集。则在 \(\mathcal{F}_{\mathfrak{S}}(X;Y)\) 的每个预紧子集 H 上，S-收敛的一致结构与在 \(B = \bigcup_{A \in \mathfrak{S}} A\) 中逐点收敛的一致结构相同。我们可以化到 \(B = X\) 且 Y 是豪斯多夫完备空间的情形；因为若 \(j\) 是典型单射 \(B \to X\) ，\(i\) 是典型映射 \(Y \to \hat{Y}\) ，则 \(\mathcal{F}(X;Y)\) 上的 S-收敛的一致结构是 \(\mathcal{F}(B;\hat{Y})\) 上的 S-收敛的一致结构在映射 \(\theta : u \to i \circ u \circ j\) 下的原像（ \(\S 1\) ，第 4 号，命题 4），且 H 预紧当且仅当 \(\theta(H)\) 预紧（第 II 章，\(\S 4\) ，第 2 号，命题 3）。在此情形，若 \(B = X\) 且 Y 是豪斯多夫完备空间，则 \(\mathcal{F}_{\mathfrak{S}}(X;Y)\) 是豪斯多夫且完备的（ \(\S 1\) ，第 2 号，命题 1 以及第 5 号，定理 1）；因此 H 在这个空间中的闭包 \(\overline{H}\) 是紧的。在 \(\overline{H}\) 上，逐点收敛拓扑是豪斯多夫的（ \(\S 1\) ，第 2 号，命题 1）且比 S-收敛拓扑粗；因此这两个拓扑相同（第 I 章，\(\S 9\) ，第 4 号，定理 2 的推论 3），从而 S-收敛与逐点收敛的一致结构也相同（第 II 章，\(\S 4\) ，第 1 号，定理 1）。

